\section{问题分析}

\subsection{总体思路}
四问的共同难点是：隐藏状态不能直接获取，观测只能逐步缩小可行集，而清除和退出又必须具有可核验的保证。因此，本文将决策分为两层：几何安全层保留所有满足观测的可能状态，并签发接收、发现、清除和退出证书；效率层只在证书允许的动作中优化观测位置和执行顺序。

问题一提供定位与清除的几何内核，问题二回答“下一次在哪里测”，问题三和问题四分别为全向源和定向源建立连续区域的发现保证。每获得一次观测，程序先更新频道可行域，再判断是否已可安全清除、是否需要补充观测，以及未知频道是否已完成覆盖审计。各问的模型衔接见\cref{tab:roadmap}。

\begin{table}[H]
  \centering
  \caption{四个问题的模型衔接与可核验输出}
  \label{tab:roadmap}
  \begin{tabularx}{\textwidth}{cLLLL}
    \toprule
    问题 & 主要输入 & 核心不确定性 & 建模方法 & 可核验输出 \\
    \midrule
    一 & 多点示向度 & 角度误差有界 & 半平面交、直径与最小包围圆 & 定位区域与安全清除证书 \\
    二 & 首次正观测 & 接收半径未知 & 保证接收约束下的主动选点 & 第二检测点与后验风险估计 \\
    三 & 多频道全向源 & 源数未知 & 七站圆盘覆盖与频道状态机 & 发现或不存在证书、清除路线 \\
    四 & 多频道定向源 & 发射方向未知 & 等边三角网格与正观测更新 & 与方向无关的连续域发现证书 \\
    \bottomrule
  \end{tabularx}
\end{table}

\subsection{问题一分析}
一次示向只限定方向而不给出距离，且误差只有确定上下界。因此不引入未知的概率分布，而将每次观测写成包含真实位置的前向角域，通过集合求交递推缩小可行域\cite{combettes1993}。定位区域为凸集时，其最远点对可由旋转卡壳计算\cite{toussaint1983}；但最远点对只回答直径问题，清除可行性仍需由最小包围圆判定。所以本问的求解链为“角域交一直径一覆盖圆”，并对空集、无界、线段等退化情形单独处理。

\subsection{问题二分析}
观测者的运动会改变仅测向定位的可观性，选点需在交会角改善与移动代价之间权衡\cite{hammel1989}。本题还有接收半径未知的约束：新位置即使几何交会角较好，也可能收不到第二次信号。因此先求对任意当前可行源位置都能再次接收的区域，再在其中联合评价移动、检测、后续处置时间和采样后验半径。有限候选采用粗到细搜索，并对入选点重新核验接收约束。

\subsection{问题三分析}
对全向源，无信号结果可根据接收半径下界排除检测点周围 $1000\,\mathrm{m}$ 圆盘。因此，连续区域搜索可先转化为有限站点覆盖问题\cite{cortes2004}：覆盖负责证明不漏检，频道状态机负责组织已发现源的补测与清除。为减少总时间，覆盖、定位和清除统一为候选宏动作。几何规则限制动作合法性，学习策略只调整合法候选的执行次序，从而保留清除和退出证书。

\subsection{问题四分析}
定向源的无信号结果既可能由距离过远造成，也可能因检测点位于发射背面而产生，所以问题三的圆盘排除不再成立。关键是建立与发射方向无关的发现证书：对目标圆域内任意源位置和任意发射方向，至少有一个检测点同时位于 $1000\,\mathrm{m}$ 内和接收半平面中。等边三角网格承担该发现任务；获得方向值后，继续使用问题一的可行域更新和安全清除方法，并沿用问题三的频道状态机。
